\section{滑模方法1}

定义滑模变量 $s_{i0}=v_{i0}+k_1p_{i0}-\varphi_i+k_2\eta_i$，其中 $\dot\eta_i=k_1p_{i0}-\varphi_i$，$k_1,k_2>0$；$\varphi_i$ 为排斥项，仅依赖相对位置。由 $\dot p_{i0}=v_{i0}$、$\dot v_{i0}=B_i(t)u_i+w_i$ 与 $\dot\eta_i$ 可得
\[
    \dot s_{i0}=B_i(t)u_i+\Delta_i+w_i,\qquad
    \Delta_i=-\dot\varphi_i+k_1v_{i0}+k_2(k_1p_{i0}-\varphi_i).
\]
逐分量写为 $\dot s_{i0,\ell}=b_{i\ell}u_{i\ell}+N_{i\ell}$（$N_{i\ell}=\Delta_{i\ell}+w_{i\ell}$）。控制律取
\[
    u_{i\ell}=R_{i\ell}\sgn\!\left[\sin\!\left(\frac{\pi}{\varepsilon}s_{i0,\ell}\right)\right],\qquad
    R_{i\ell}=\frac{|\Delta_{i\ell}|+\bar w_i+\lambda}{\ub_{i\ell}},
\]
其中 $\lambda>0$。该律不使用 $b_{i\ell}$ 的符号信息。

\subsection{到达性与有界性}

固定 $i,\ell$，记 $s=s_{i0,\ell}$、$b=b_{i\ell}$、$N=\Delta+w$。闭环为 $\dot s=bR\sgn[\sin(\pi s/\varepsilon)]+N$。由 $R=(|\Delta|+\bar w_i+\lambda)/\ub$ 与 $|b|\ge\ub$ 得 $|b|R\ge|\Delta|+\bar w_i+\lambda$，而 $|N|\le|\Delta|+\bar w_i$，故 $|b|R-|N|\ge\lambda$，即 $|b|R>|N|$。

$b>0$ 时奇数倍面 $s=(2r+1)\varepsilon$ 吸引；$b<0$ 时偶数倍面 $s=2r\varepsilon$ 吸引。取 $\tilde s=s-k\varepsilon$ 且 $|\tilde s|<\varepsilon$，两种情况均统一化为 $\dot{\tilde s}=-|b|R\,\sgn(\tilde s)+N$。取 $V=\tfrac12\tilde s^2$ 得 $\dot V\le-(|b|R-|N|)|\tilde s|\le-\lambda|\tilde s|=-\lambda\sqrt{2V}$，故有限时间到达，且到达时间 $T_r\le|\tilde s(t_0)|/\lambda$。由此可知 $s_{i0}\in L_\infty$。Filippov 等效控制 $u_{\mathrm{eq}}=-N/b$ 满足 $|u_{\mathrm{eq}}|<R$，故滑模面上 $\dot s_{i0}=0$。若 $b_{i\ell}$ 有限次切换则每次重新到达，最后一次切换后进入周期滑模；无限切换需额外驻留时间假设。

\subsection{合围性质证明}

排斥项取 $\varphi_i=\sum_{j\in\mathcal N_i}\rho(\|p_{ij}\|)p_{ij}/\|p_{ij}\|$，其中 $\rho(s)=\max\{1/(s-d)-1/(\mu-d),0\}$；无向图下成对抵消，故 $\sum_i\varphi_i=0$。

\textbf{瞬态避碰（P2）。} 定义增广势函数
\[
    V_W=\frac{k_1}{2}\sum_i\|p_{i0}\|^2
    +\frac12\sum_i\sum_{j\in\mathcal N_i}\int_{\|p_{ij}\|}^{\mu}\rho\,ds
    +\frac{k_2}{2}\sum_i\|\eta_i\|^2 .
\]
由无向图对称性，其导数满足 $\dot V_W=\sum_i g_i^Tv_{i0}+k_2\sum_i\eta_i^T\dot\eta_i$，其中 $g_i=k_1p_{i0}-\varphi_i=\dot\eta_i$。代入 $v_{i0}=s_{i0}-g_i-k_2\eta_i$ 得
\[
    \dot V_W=-\sum_i\|g_i\|^2+\sum_i g_i^Ts_{i0}
    \le\frac14\sum_i\|s_{i0}\|^2 .
\]
因 $s_{i0}\in L_\infty$，故 $\dot V_W\le C$。于是 $V_W$ 在任意有限区间有界；而 $\|p_{ij}\|\to d^+$ 时势垒项趋于无穷，矛盾，故闭环不碰撞。

\textbf{稳态合围（P1）。} 滑模面上有 $v_{i0}+k_1p_{i0}-\varphi_i+k_2\eta_i=k_i\varepsilon_i$ 为常值。对全体平均得 $\dot{\tilde p}+k_1\tilde p+k_2\bar\eta=\bar s$、$\dot{\bar\eta}=k_1\tilde p$（$\tilde p=\bar p-p_0$，$\bar\eta=\frac1N\sum\eta_i$），故 $\ddot{\tilde p}+k_1\dot{\tilde p}+k_1k_2\tilde p=0$。其特征多项式 $\lambda^2+k_1\lambda+k_1k_2$ 在 $k_1,k_2>0$ 时为 Hurwitz，于是 $\tilde p\to0$，即质心趋近目标；质心位于凸包内，故 $\mathrm{dist}(p_0,\mathrm{co}(p))\to0$。

说明：本方法滑模变量直接含 $-\varphi_i$，求导后出现 $\dot\varphi_i$，需邻居相对速度，与“仅测相对位置”假设矛盾；且 P3（速度一致）未证。
